\section{价格事实：预告日涨停后，市场迅速回到持续性问题}

公司在 7 月 6 日晚发布业绩预告后，股票在 7 月 7 日收于 24.65 元、单日涨幅 10.00\%。此后价格回撤并伴随高换手：7 月 8 日下跌 9.98\%，7 月 9 日下跌 8.97\%，至 7 月 22 日收于 16.79 元，较 7 月 7 日回撤 31.9\%。价格路径本身不能证明资金动机或基本面因果；它使 H2 持续性、估值和现金质量成为正式 H1 需要优先验证的假说，而不是已确认的市场动机。

\begin{exhibitbox}[业绩预告后的可复核价格路径]
\begin{tabularx}{\exhibitboxwidth}{L{2.4cm}R{2.1cm}R{2.0cm}R{2.3cm}X}
\toprule
交易日 & 收盘价 & 当日涨跌 & 换手率 & 读法 \\
\midrule
7月6日 & 22.41 元 & +2.47\% & 5.44\% & 公告前一日 \\
7月7日 & 24.65 元 & +10.00\% & 1.67\% & 预告后涨停 \\
7月8日 & 22.19 元 & -9.98\% & 12.25\% & 高换手回吐开始 \\
7月9日 & 20.20 元 & -8.97\% & 11.60\% & 价格继续下修 \\
7月22日 & 16.79 元 & -0.71\% & 3.86\% & 较 7 月 7 日 -31.9\% \\
\bottomrule
\end{tabularx}
\end{exhibitbox}
\sourcenote{东财 002497 日线快照，2026-07-23 抓取；收益率按未复权收盘价计算。}

最近 20 个交易日（2026-06-25 至 2026-07-22）最高/最低价为 24.65/15.87 元，日均成交额约 13.63 亿元、日均换手率约 6.47\%。这是一只围绕业绩预期高度交易的股票，而不是可以忽略事件风险的低波动标的。

\section{市场在交易什么：证据、推断与未知}

\begin{tabularx}{\linewidth}{L{2.7cm}X X}
\toprule
层级 & 内容 & 对行动的意义 \\
\midrule
已确认事实 & H1 预告利润跳升；行业锂价/需求环境好于 2025 年末；同业 Q1 也受益于锂价 & Q2 盈利弹性不是无根据叙事 \\
研究推断 & 预告日后快速回撤与高换手，\textbf{可能}反映市场将重心由 H1 惊喜转向 H2 持续性、估值和现金质量 & 需要用正式 H1 数据验证，不能当作已证实因果 \\
尚未知道 & H2 实际销量、ASP、单位成本、自给率、套保、民爆订单及现金回收 & 不能把低 PE 或高 Q2 利润直接翻译为买入 \\
\bottomrule
\end{tabularx}

\section{外部观点：不是“共识”，而是三种 H2 要求}

三家完整券商报告的全年归母预测为 23.15--30.32 亿元；东吴是唯一给出 42 元目标价的机构，方法为其自身 16 倍 2026E PE，报告日价格为 24.65 元。这个目标价既不是一致预期，也不是 7 月 22 日的当前上行空间。对事件研究更有价值的是三家隐含 H2 利润区间的差异：它把“利好预告”具体化为需要通过正式 H1 数据验证的销量、价格、成本、资源与民爆假设。
